\chapter{Elementos Textuais}

De acordo com a ABNT NBR 14724:2024, o texto é composto por uma parte introdutória, pelo desenvolvimento e por uma parte conclusiva. A norma deixa a nomenclatura dos títulos a critério do autor e orienta a organização do trabalho em seções. Portanto, a sequência apresentada neste capítulo é um modelo acadêmico recomendável, não uma lista de títulos obrigatórios para todas as áreas.

\section{Introdução}

A introdução deve situar o leitor. Ela apresenta tema, delimitação, problema, objetivos, justificativa, síntese do método e organização do trabalho.

Uma introdução eficiente responde, de modo articulado, às seguintes perguntas:

\begin{itemize}
  \item Qual é o tema e qual recorte foi adotado?
  \item Qual problema será enfrentado?
  \item Por que o problema é relevante?
  \item O que o trabalho pretende alcançar?
  \item Como a investigação foi conduzida?
  \item Como o texto está organizado?
\end{itemize}

A introdução abre o percurso e prepara o leitor para as seções seguintes. Ela não deve antecipar toda a discussão nem se transformar em revisão bibliográfica completa.

\section{Desenvolvimento}

O desenvolvimento é a parte central do trabalho. Sua divisão depende da natureza da pesquisa, da modalidade do TCC e das práticas da área. As seções a seguir representam uma organização frequente e podem ser adaptadas com o orientador.

\subsection{Referencial teórico ou revisão de literatura}

O referencial apresenta conceitos, autores e debates que sustentam a análise. Deve ser seletivo e articulado ao problema. Não basta resumir textos em sequência; é necessário construir diálogo entre as fontes, explicitar divergências e mostrar como o trabalho se posiciona.

\subsection{Metodologia}

A metodologia descreve como o trabalho foi realizado e deve permitir que o leitor avalie a adequação do percurso. Conforme o tipo de pesquisa, pode indicar:

\begin{itemize}
  \item abordagem e desenho da pesquisa;
  \item população, amostra, corpus, documentos, bases ou materiais;
  \item critérios de inclusão e exclusão;
  \item instrumentos e procedimentos de coleta;
  \item procedimentos de tratamento e análise;
  \item cuidados éticos e proteção de dados;
  \item limitações metodológicas.
\end{itemize}

\caixaorientacao{A metodologia deve ser escrita no nível de detalhe suficiente para que outro leitor compreenda o caminho percorrido. Frases genéricas como ``foi feita pesquisa na internet'' não atendem à exigência acadêmica.}

\subsection{Dados ou corpus: seção opcional}

Uma seção primária própria para dados ou corpus pode ser inserida depois da metodologia quando o conjunto analisado exigir caracterização extensa. Ela pode descrever origem, variáveis, composição da amostra, critérios de qualidade, dados ausentes, transformações, organização do corpus e produto final do tratamento.

Essa seção é opcional. Se a caracterização for breve, deve permanecer como subseção da metodologia ou ser incorporada aos resultados. Separá-la só se justifica quando melhora a compreensão e evita que a metodologia fique desproporcional. Ela não deve repetir o procedimento de coleta nem antecipar a interpretação dos achados.

\subsection{Resultados e discussão}

Os resultados apresentam o que foi encontrado, produzido ou sistematizado. A discussão interpreta esses achados à luz do problema, dos objetivos e da literatura, explicando significados, convergências, diferenças, implicações e limites.

Há duas organizações igualmente possíveis:

\begin{itemize}
  \item \textbf{Resultados e discussão integrados:} cada achado é apresentado e imediatamente interpretado. Essa forma costuma favorecer estudos qualitativos, aplicados ou organizados por temas.
  \item \textbf{Resultados e discussão separados:} primeiro são apresentados os achados e depois é feita a interpretação conjunta. Essa forma pode ser útil quando há grande volume de dados ou tradição disciplinar específica.
\end{itemize}

A escolha deve ser feita com o orientador. Resultados não devem ser acumulados sem explicação, e a discussão não deve introduzir dados que não foram apresentados.

\section{Conclusão ou considerações finais}

A parte conclusiva retoma o problema e os objetivos, sintetiza os principais achados e explicita a resposta que o trabalho alcançou. Também pode registrar contribuições, aplicações, limitações e recomendações ou estudos futuros coerentes com os resultados.

Os títulos ``Conclusão'' e ``Considerações finais'' são admissíveis quando compatíveis com a orientação do curso. Em geral, ``conclusão'' comunica um fechamento mais direto, enquanto ``considerações finais'' pode ser apropriado quando o objeto admite síntese reflexiva sem pretensão de encerramento definitivo.

Essa parte não deve apresentar teoria nova, dados inéditos ou discussão que não tenha sido desenvolvida anteriormente. Um bom fechamento responde ao que o trabalho prometeu investigar, sem repetir mecanicamente o resumo ou a introdução.
